% !TEX root = ../../../main.tex
% 所属：第四部分：Fourier 分析 / 第6章：有限Fourier分析
% 节：6.1 $\Z(N)$上的Fourier分析
% 正文按原文件顺序保存；由 ../chapter6.tex 调用。
\section{$\Z(N)$上的Fourier分析}

\begin{definition}\label{def:finite-cyclic-group-ZN}
	所有 $N$ 次单位根组成的集合记为
	\[
	\Z(N)
	=\left\{1,e^{2\pi i/N},\ldots,e^{2\pi i(N-1)/N}\right\}.
	\]
	在复数乘法下，$\Z(N)$ 是一个 Abelian 群，并且与
	$\Z/N\Z$ 同构。
\end{definition}

记
\[
\zeta_N=e^{2\pi i/N},
\]
并设 $V$ 为 $\Z(N)$ 上所有复值函数组成的向量空间。以后简记
\[
F(k)=F(\zeta_N^k),
\qquad 0\leq k\leq N-1.
\]
函数 $F\in V$ 由
\[
F(0),F(1),\ldots,F(N-1)
\]
唯一确定，因此 $V\cong\C^N$，且 $\dim_\C V=N$。

在 $V$ 上定义 Hermitian 内积与范数
\[
\langle F,G\rangle
=\sum_{k=0}^{N-1}F(k)\overline{G(k)},
\qquad
\norm{F}^2=\sum_{k=0}^{N-1}\abs{F(k)}^2.
\]

类似于 $\R$ 上的 Fourier 分析，下面寻找 $V$ 的一组标准正交基。
对 $0\leq\ell\leq N-1$，定义
\[
e_\ell(k)=\zeta_N^{\ell k}
=e^{2\pi i\ell k/N},
\qquad 0\leq k\leq N-1.
\]

\begin{lemma}\label{lem:finite-fourier-orthogonality}
	函数组 $\{e_0,\ldots,e_{N-1}\}$ 两两正交。更准确地，
	\[
	\langle e_m,e_\ell\rangle
	=\begin{cases}
	N,&m=\ell,\\
	0,&m\neq\ell.
	\end{cases}
	\]
\end{lemma}

\begin{proof}
	由定义，
	\[
	\langle e_m,e_\ell\rangle
	=\sum_{k=0}^{N-1}\zeta_N^{mk}\zeta_N^{-\ell k}
	=\sum_{k=0}^{N-1}\zeta_N^{(m-\ell)k}.
	\]
	当 $m=\ell$ 时，每一项均为 $1$，故和为 $N$。当
	$m\neq\ell$ 时，令 $q=\zeta_N^{m-\ell}$，则 $q\neq1$ 且
	$q^N=1$，从而
	\[
	1+q+\cdots+q^{N-1}=\frac{1-q^N}{1-q}=0.
	\]
\end{proof}

由于 $\dim_\C V=N$，引理 \ref{lem:finite-fourier-orthogonality}
表明 $\{e_0,\ldots,e_{N-1}\}$ 是 $V$ 的一组正交基。令
\[
e_\ell^*=\frac1{\sqrt N}e_\ell,
\]
则 $\{e_0^*,\ldots,e_{N-1}^*\}$ 是一组标准正交基。因此对任意
$F\in V$，
\[
F=\sum_{n=0}^{N-1}\langle F,e_n^*\rangle e_n^*,
\qquad
\norm{F}^2
=\sum_{n=0}^{N-1}\abs{\langle F,e_n^*\rangle}^2.
\]

\begin{definition}\label{def:finite-fourier-coefficients}
	函数 $F\in V$ 的第 $n$ 个 Fourier 系数定义为
	\[
	a_n
	=\frac1N\sum_{k=0}^{N-1}
	F(k)e^{-2\pi ikn/N},
	\qquad 0\leq n\leq N-1.
	\]
\end{definition}

\begin{theorem}(Fourier 反演与 Parseval--Plancherel 公式)
	\label{thm:finite-fourier-inversion-parseval}
	设 $F$ 是 $\Z(N)$ 上的函数，则
	\[
	F(k)=\sum_{n=0}^{N-1}a_ne^{2\pi ink/N},
	\qquad 0\leq k\leq N-1.
	\]
	此外，
	\[
	\sum_{n=0}^{N-1}\abs{a_n}^2
	=\frac1N\sum_{k=0}^{N-1}\abs{F(k)}^2.
	\]
\end{theorem}

\begin{proof}
	由定义，
	\[
	a_n=\frac1N\langle F,e_n\rangle
	=\frac1{\sqrt N}\langle F,e_n^*\rangle.
	\]
	将其代入标准正交基展开式与 Parseval 恒等式，立即得到上述两个公式。
\end{proof}


	圆周上充分光滑的函数，例如 $C^2$ 函数，其 Fourier 逆变换
	可以通过令模 $N\to\infty$，把上述有限和看作 Riemann 和得到。具体有以下的命题。





令 $f$ 是定义在单位圆周 $S^1$ 上的函数。对每个 $N\geq1$，定义
有限 Fourier 系数
\[
a_N(n)=\frac1N\sum_{k=1}^N
f\left(e^{2\pi ik/N}\right)e^{-2\pi ikn/N},
\qquad n\in\Z.
\]
同时记
\[
a(n)=\int_0^1 f\left(e^{2\pi ix}\right)e^{-2\pi inx}\dd x
\]
为 $f$ 通常的 Fourier 系数。

\begin{proposition}\label{prop:finite-circle-fourier-limit}
	上述有限 Fourier 系数满足
	\[
	a_N(n+N)=a_N(n),
	\qquad n\in\Z.
	\]
	若 $f$ 连续，则对每个固定的 $n\in\Z$，当 $N\to\infty$ 时，
	\[
	a_N(n)\longrightarrow a(n).
	\]
0\end{proposition}

\begin{proof}
	首先，
	\[
	\begin{aligned}
	a_N(n+N)
	&=\frac1N\sum_{k=1}^N
	f\left(e^{2\pi ik/N}\right)e^{-2\pi ik(n+N)/N}\\
	&=\frac1N\sum_{k=1}^N
	f\left(e^{2\pi ik/N}\right)e^{-2\pi ikn/N}
	=a_N(n).
	\end{aligned}
	\]
	若 $f$ 连续，则
	\[
	g_n(x)=f\left(e^{2\pi ix}\right)e^{-2\pi inx}
	\]
	在 $[0,1]$ 上连续，而
	\[
	a_N(n)=\frac1N\sum_{k=1}^Ng_n\left(\frac{k}{N}\right)
	\]
	是 $g_n$ 的 Riemann 和，故它收敛到
	$\displaystyle\int_0^1g_n(x)\dd x=a(n)$。
\end{proof}

\begin{proposition}\label{prop:finite-circle-fourier-C1-decay}
	若 $f\in C^1(S^1)$，则存在只依赖于 $f$ 的常数 $c>0$，使得
	\[
	\abs{a_N(n)}\leq\frac{c}{\abs n},
	\qquad 0<\abs n\leq\frac N2.
	\]
\end{proposition}

\begin{proof}
	令
	\[
	F(x)=f\left(e^{2\pi ix}\right),
	\qquad
	M_1=\sup_{x\in\R}\abs{F'(x)}.
	\]
	函数 $F$ 是以 $1$ 为周期的 $C^1$ 函数。取整数 $\ell$ 使得
	\[
	\abs{\ell-\frac{N}{2\abs n}}\leq\frac12.
	\]
	由 $0<\abs n\leq N/2$ 可得
	\[
	\frac14\leq\frac{\ell\abs n}{N}\leq\frac34.
	\]
	利用 $F$ 的周期性，将有限和的指标平移 $\ell$，得到
	\[
	\begin{aligned}
	a_N(n)\left(1-e^{2\pi i\ell n/N}\right)
	=\frac1N\sum_{k=1}^N
	\left[F\left(\frac{k}{N}\right)
	-F\left(\frac{k+\ell}{N}\right)\right]
	e^{-2\pi ikn/N}.
	\end{aligned}
	\]
	一方面，由中值定理，
	\[
	\abs{F\left(\frac{k}{N}\right)
	-F\left(\frac{k+\ell}{N}\right)}
	\leq M_1\frac{\ell}{N}.
	\]
	另一方面，
	\[
	\abs{1-e^{2\pi i\ell n/N}}
	=2\abs{\sin\frac{\pi\ell n}{N}}
	\geq\sqrt2.
	\]
	因此
	\[
	\abs{a_N(n)}
	\leq\frac{M_1}{\sqrt2}\frac{\ell}{N}
	\leq\frac{3M_1}{4\sqrt2}\frac1{\abs n},
	\]
	取 $c=3M_1/(4\sqrt2)$ 即可。
\end{proof}

\begin{proposition}\label{prop:finite-circle-fourier-C2-decay-inversion}
	若 $f\in C^2(S^1)$，则存在只依赖于 $f$ 的常数 $c>0$，使得
	\[
	\abs{a_N(n)}\leq\frac{c}{\abs n^2},
	\qquad 0<\abs n\leq\frac N2.
	\]
	此外，$f$ 满足 Fourier 反演公式
	\[
	f\left(e^{2\pi ix}\right)
	=\sum_{n=-\infty}^{\infty}a(n)e^{2\pi inx}.
	\]
\end{proposition}

\begin{proof}
	仍令 $F(x)=f(e^{2\pi ix})$，并记
	\[
	M_2=\sup_{x\in\R}\abs{F''(x)}.
	\]
	沿用上一个证明中 $\ell$ 的选取，令
	\[
	h=\frac{\ell}{N},
	\qquad
	\theta=\frac{2\pi\ell n}{N}.
	\]
	将求和指标分别向两侧平移 $\ell$，可得
	\[
	\left(e^{i\theta}+e^{-i\theta}-2\right)a_N(n)
	=\frac1N\sum_{k=1}^N
	\left[F\left(\frac{k}{N}+h\right)
	+F\left(\frac{k}{N}-h\right)
	-2F\left(\frac{k}{N}\right)\right]
	e^{-2\pi ikn/N}.
	\]
	Taylor 公式给出
	\[
	\abs{F(x+h)+F(x-h)-2F(x)}\leq M_2h^2.
	\]
	又因为
	\[
	\abs{e^{i\theta}+e^{-i\theta}-2}
	=4\sin^2\frac{\theta}{2}
	=4\sin^2\frac{\pi\ell n}{N}
	\geq2,
	\]
	所以
	\[
	\abs{a_N(n)}
	\leq\frac{M_2}{2}h^2
	\leq\frac{9M_2}{32}\frac1{\abs n^2}.
	\]

	下面由有限 Fourier 反演公式推出圆周上的反演公式。只取奇数
	$N$，并对给定的 $x\in[0,1]$ 选取整数 $k_N$，使得
	$k_N/N\to x$。有限 Fourier 反演公式和
	命题 \ref{prop:finite-circle-fourier-limit} 的周期性给出
	\[
	F\left(\frac{k_N}{N}\right)
	=\sum_{\abs n\leq(N-1)/2}
	a_N(n)e^{2\pi ink_N/N}.
	\]
	固定 $M\in\Z^+$。当 $N\to\infty$ 时，有限项部分满足
	\[
	\sum_{\abs n\leq M}a_N(n)e^{2\pi ink_N/N}
	\longrightarrow
	\sum_{\abs n\leq M}a(n)e^{2\pi inx}.
	\]
	而由上面的二阶衰减估计，
	\[
	\abs{\sum_{M<\abs n\leq(N-1)/2}
	a_N(n)e^{2\pi ink_N/N}}
	\leq\sum_{\abs n>M}\frac{c}{n^2},
	\]
	右端在 $M\to\infty$ 时趋于 $0$。又因 $F$ 连续，
	$F(k_N/N)\to F(x)$。先令 $N\to\infty$，再令
	$M\to\infty$，便得到
	\[
	f\left(e^{2\pi ix}\right)=F(x)
	=\sum_{n=-\infty}^{\infty}a(n)e^{2\pi inx}.
	\]
\end{proof}
也就是说，连续的Fourier级数可以看成有限Fourier分析的极限。
\subsection*{快速 Fourier 变换}

快速 Fourier 变换是在计算 $\Z(N)$ 上函数的 Fourier 系数时
产生和发展起来的，它已经成为数值分析中的一个重要工具。
下面估计计算一组 Fourier 系数所需的运算次数。

固定 $N$，设 $F(0),\ldots,F(N-1)$ 已知，并记
\[
\omega_N=e^{-2\pi i/N}.
\]
函数 $F$ 的第 $k$ 个 Fourier 系数写为
\[
a_k^N(F)
=\frac1N\sum_{r=0}^{N-1}F(r)\omega_N^{kr}.
\]
记 $\#(N)$ 为计算全部 $N$ 个 Fourier 系数所需的运算次数。
直接计算时，需要预先求出
\[
\omega_N^2,\ldots,\omega_N^{N-1},
\]
再逐项求和。粗略估计有
\[
\#(N)<2N^2+N,
\]
因而直接算法的复杂度为 $O(N^2)$。

若 $N=2^n$，可以利用偶数项和奇数项的分解显著减少运算次数。

\begin{theorem}[快速 Fourier 变换的复杂度]
	\label{thm:fft-complexity}
	若 $N=2^n$，其中 $n\geq1$，则
	\[
	\#(N)\leq4N\log_2N=4Nn.
	\]
\end{theorem}

\begin{lemma}\label{lem:fft-recursion}
	设 $\omega_{2M}=e^{-2\pi i/(2M)}$，则
	\[
	\#(2M)\leq2\#(M)+8M.
	\]
\end{lemma}

\begin{proof}
	计算 $\omega_{2M},\ldots,\omega_{2M}^{2M}$ 至多需要
	$2M$ 次运算。注意
	\[
	\omega_M=e^{-2\pi i/M}=\omega_{2M}^2.
	\]
	对给定的 $F:\Z(2M)\to\C$，定义两个 $\Z(M)$ 上的函数
	\[
	F_0(r)=F(2r),
	\qquad
	F_1(r)=F(2r+1),
	\qquad 0\leq r\leq M-1.
	\]
	将 Fourier 系数中的求和分成偶数项与奇数项，得到
	\[
	\begin{aligned}
	a_k^{2M}(F)
	&=\frac1{2M}\sum_{r=0}^{2M-1}F(r)\omega_{2M}^{kr}\\
	&=\frac12\left[
	\frac1M\sum_{\ell=0}^{M-1}F(2\ell)\omega_{2M}^{2k\ell}
	+\frac1M\sum_{m=0}^{M-1}
	F(2m+1)\omega_{2M}^{k(2m+1)}
	\right]\\
	&=\frac12\left[
	\frac1M\sum_{\ell=0}^{M-1}F_0(\ell)\omega_M^{k\ell}
	+\frac1M\sum_{m=0}^{M-1}
	F_1(m)\omega_M^{km}\omega_{2M}^k
	\right]\\
	&=\frac12\left[a_k^M(F_0)+a_k^M(F_1)\omega_{2M}^k\right].
	\end{aligned}
	\]
	这里 $a_k^M$ 的下标按模 $M$ 理解。计算 $F_0,F_1$ 的全部
	Fourier 系数需要 $2\#(M)$ 次运算；得到每个
	$a_k^{2M}(F)$ 至多再需三次运算。因此
	\[
	\#(2M)\leq2M+2\#(M)+3(2M)
	=2\#(M)+8M.
	\]
\end{proof}

\begin{proof}[定理的证明]
	对 $n$ 归纳。当 $n=1$ 时，$N=2$，两个 Fourier 系数为
	\[
	a_0^2(F)=\frac12\bigl(F(1)+F(-1)\bigr),
	\qquad
	a_1^2(F)=\frac12\bigl(F(1)-F(-1)\bigr).
	\]
	计算它们至多需要 $5$ 次运算，故 $\#(2)\leq8$。

	假设结论对 $2^{n-1}$ 成立。令 $M=2^{n-1}$，由引理
	\ref{lem:fft-recursion}，
	\[
	\begin{aligned}
	\#(2^n)
	&\leq2\#(2^{n-1})+8\cdot2^{n-1}\\
	&\leq2\cdot4\cdot2^{n-1}(n-1)+8\cdot2^{n-1}\\
	&=4\cdot2^n n.
	\end{aligned}
	\]
	归纳完成。
\end{proof}

